% Shared packages / theorem environments for standalone topic articles.
% Body content must live in each topic's .tex (do not \input chapters/).

\usepackage[a4paper,margin=2.35cm]{geometry}
\usepackage{amsmath,amssymb,amsthm,mathtools}
\usepackage{bm}
\usepackage{enumitem}
\usepackage{booktabs}
\usepackage{longtable}
\usepackage{array}
\usepackage{xcolor}
\usepackage{tikz}
\usetikzlibrary{trees,positioning,shapes.geometric,calc}
\usepackage{xurl}
\usepackage{hyperref}
\usepackage{csquotes}
\usepackage[backend=biber,style=authoryear,sorting=nyt,maxcitenames=2,maxbibnames=6]{biblatex}

\addbibresource{refs.bib}

\hypersetup{
  colorlinks=true,
  linkcolor=blue!50!black,
  citecolor=blue!50!black,
  urlcolor=blue!60!black
}

\setlist[itemize]{topsep=3pt,itemsep=2pt,parsep=0pt}
\setlist[enumerate]{topsep=3pt,itemsep=2pt,parsep=0pt}
\setcounter{tocdepth}{2}
\setlength{\emergencystretch}{3em}
\raggedbottom
\newcolumntype{L}[1]{>{\raggedright\arraybackslash}p{#1}}

% --- standard amsthm environments (section-numbered) ---
\theoremstyle{plain}
\newtheorem{theorem}{定理}[section]
\newtheorem{lemma}[theorem]{引理}
\newtheorem{proposition}[theorem]{命题}
\newtheorem{corollary}[theorem]{推论}

\theoremstyle{definition}
\newtheorem{definition}{定义}[section]
\newtheorem{example}[definition]{例}
\newtheorem{teaching}[definition]{从零教学}

\theoremstyle{remark}
\newtheorem{remark}{注}[section]
\newtheorem{heuristic}[remark]{启发式解释}
\newtheorem{engineering}[remark]{工程惯例}
\newtheorem{examnote}[remark]{考试提醒}
\newtheorem{checkpoint}[remark]{考点}
\newtheorem{proofsketch}[remark]{证明骨架}

% Keep examproblem available if a topic ever needs it
\newtheoremstyle{examstyle}
  {6pt}{6pt}{\normalfont}{}{\bfseries}{.}{0.5em}{}
\theoremstyle{examstyle}
\newtheorem{examproblem}{题}[section]

% --- math macros ---
\newcommand{\R}{\mathbb{R}}
\newcommand{\E}{\mathbb{E}}
\newcommand{\Pbb}{\mathbb{P}}
\newcommand{\N}{\mathcal{N}}
\newcommand{\Hcal}{\mathcal{H}}
\newcommand{\Dcal}{\mathcal{D}}
\newcommand{\Lcal}{\mathcal{L}}
\newcommand{\argmin}{\operatorname*{arg\,min}}
\newcommand{\argmax}{\operatorname*{arg\,max}}
\newcommand{\prox}{\operatorname{prox}}
\newcommand{\Var}{\operatorname{Var}}
\newcommand{\KL}{\operatorname{KL}}
\newcommand{\softmax}{\operatorname{softmax}}
\newcommand{\Dir}{\operatorname{Dir}}
\newcommand{\Cat}{\operatorname{Cat}}
\newcommand{\Bern}{\operatorname{Bern}}
\newcommand{\tr}{\operatorname{tr}}

\newcommand{\source}[1]{\par\smallskip\noindent{\small\textit{来源定位：#1}}\par\smallskip}
\newcommand{\solution}{\par\smallskip\noindent\textbf{解答。}\ }
\newcommand{\officialenglish}[1]{\par\smallskip\noindent\textbf{Official English prompt.}\ \begin{quote}\small #1\end{quote}}
\newcommand{\chinesetranslation}[1]{\par\smallskip\noindent\textbf{中文翻译。}\ #1\par\smallskip}
\newcommand{\concepttarget}[1]{\phantomsection\label{#1}\hypertarget{#1}{}}
\newcommand{\conceptref}[2]{\hyperref[#1]{#2}}
